\documentclass[../main.tex]{subfiles}

\graphicspath{{\subfix{../images/}}}

\begin{document}

\clearpage
\thispagestyle{empty}
\section{Diagrammi di Bode}

Procediamo ora con la seconda parte del corso andando ad introdurre alcune tipologie di digrammi utili allo studio della materia.


\subsection{Introduzione}

Prima di approfondire gli argomenti trattati in questa sezione, chiariamo esattamente di cosa parleremo. In particolare, vogliamo fornire una spiegazione chiara e accurata, seppur breve, di come funzionano i diagrammi di Bode, cosa rappresentano e come vengono utilizzati in questo corso.\\


I diagrammi di Bode sono una rappresentazione grafica in scala logaritmica della risposta in frequenza di un sistema dinamico lineare e tempo-invariante (LTI). Si usano due grafici separati: modulo (in dB) e fase (in gradi) in funzione della frequenza angolare $\omega$ (rad/s) su scala logaritmica. In automatica sono fondamentali perché permettono di analizzare stabilità, prestazioni e robustezza di un sistema di controllo in modo intuitivo e quantitativo, senza dover risolvere direttamente equazioni differenziali nel dominio del tempo.


\subsection{Risposta in frequenza}

\subsection{Diagrammi di Bode e FdT elementari}


\subsubsection{DdB di un guadagno}
\subsubsection{DdB di un polo nell'origine}
\subsubsection{DdB di un polo reale}
\subsubsection{DdB di poli complessi coniugati}

\subsection{Osservazioni}

Arrivati a questo punto andiamo a fare dell osservazioni su quanto abbiamo appena detto in questa sezione.
Come prima cosa, dopo aver visto i vari DdB di FdT elementari specifichiamo che:


\begin{enumerate}
	\item[\textbf{-}] I DdB degli zeri si ottengono da quelli dei poli di pari valore cambiando segno sia ai moduli espressi in dB che alle fasi espresse in gradi.
	
	\item[\textbf{-}] I DdB dei poli (zeri) di molteplicità \textbf{i} si ottengono da quelli del polo (zero) semplice moltiplicando i valori del modulo e della fase per \textbf{i}.
	
	
	\item[\textbf{-}] I DdB del modulo tende a $\infty$ dB ($\infty& \unit{u_{nat}})$ $\Leftrightarrow$ c'è almeno un polo sull'asse immaginario del piano complesso.
		\begin{enumerate}
			\item[\textbf{-}] Polo reale in zero, semplice o multiplo 
			\item[\textbf{-}] Coppia di poli immaginari coniugati $\pm j\omega_$  
		\end{enumerate}
	
\end{enumerate}

Infine considerando i DdB di una generica FdT $G(j\omega)$ possiamo dire che:

\begin{enumerate}
	\item[\textbf{-}] Può essere espressa come prodotto di FdT elementari dette fattori.
	
	\item[\textbf{-}] Il \textbf{modulo} di $G$ in dB in una qualunque $\omega$ é dato dalla \textbf{somma dei moduli in dB} delle FdT elementari nella stessa $\omega$.
	
	\item[\textbf{-}] La \textbf{fase} di $G$ in una qualunque $\omega$ è data dalla \textbf{somma delle fasi} delle FdT, indicate come $G_i$, nella stessa $\omega$.
	
	\item[\textbf{-}] I diagrammi di Bode di una generica FdT $G$ possono essere quindi essere costruiti sommando i contributi delle singole FdT elementari.
	
\end{enumerate}

Infine, come vedremo più avanti, i diagrammi di Bode possono essere tracciati in ambiante Matlab utilizzando il relativo comando riportato di seguito.\\

Il comando \texttt{bode} in ambiente Matlab calcola e disegna il diagramma di Bode (modulo in dB e fase in gradi) della funzione di trasferimento o del modello dinamico che gli viene passato.\\


\lstinputlisting[language = Matlab, firstline = 3, lastline = 3]{code/ex.m}
\hspace{0.5cm}


Il comando appena riportato genera il grafico di modulo e fase della FdT da noi passata come parametro. Per fare un esempio andiamo a generare una funzione di trasferimento $G$ così definita:

\[ G = \frac{10}{s \cdot (s+5) }\]\

In Matlab si ottiene che:\\

\lstinputlisting[language = Matlab, firstline = 7, lastline = 9]{code/ex.m}
\hspace{0.5cm}

Infine ricordiamo che la funzione \texttt{ft} é utilizzabile dopo l'installazione del pacchetto \textbf{Control System Toolbox}, inoltre alla funzione vengono passati come parametri prima il numeratore e poi il denominatore della FdT che volgiamo generare. Facciamo notare infine come sia anche possibile ottenere i valori numerici ritornati dal comando \texttt{bode} nel modo seguente:\\

\lstinputlisting[language = Matlab, firstline = 9, lastline = 11]{code/ex.m}
\hspace{0.5cm}


Se tutto é stato fatto nel modo corretto si dovrebbero visualizzare i seguenti grafici riportati di seguito.

 


\clearpage
\thispagestyle{empty}
\section{Diagrammi di polari e di Nyquist}








\clearpage
\thispagestyle{empty}
\section{Criterio di Nyquist}





\end{document}